\subsection{EXISTENCE THEOREM}

We retain the hypotheses and notation of the preceding number. Recall that if \(\alpha, \beta \in \mathrm{B}\) and if \(\alpha \neq \beta\) , then \(n(\beta, \alpha) \leq 0\) ; moreover, if \(n(\beta, \alpha) = 0\) , then \(n(\alpha, \beta) = 0\) (Chap. VI, §1, no. 1, formula (8)). For any pair \((\alpha, \beta)\) of distinct elements of B, put

\[
x _ {\alpha \beta} = (\mathrm{ad} x _ {\alpha}) ^ {1 - n (\beta , \alpha)} x _ {\beta} y _ {\alpha \beta} = (\mathrm{ad} x _ {- \alpha}) ^ {1 - n (\beta , \alpha)} x _ {- \beta}.
\]

Then \(x_{\alpha \beta} \in \mathfrak{a}_{+}, y_{\alpha \beta} \in \mathfrak{a}_{-}\) . If \(\theta\) denotes the canonical automorphism of \(\mathfrak{a}\) , \(\theta(x_{\alpha \beta}) = y_{\alpha \beta}\) .

Lemma 6. Let \(\alpha, \beta \in \mathrm{B}\) with \(\alpha \neq \beta\) . Then

\[
[ \mathfrak {a} _ {+}, y _ {\alpha \beta} ] = 0 \quad [ \mathfrak {a} _ {-}, x _ {\alpha \beta} ] = 0.
\]

The second formula follows from the first by using the automorphism \(\theta\) . To prove the first, it suffices to show that \([x_{\gamma}, y_{\alpha \beta}] = 0\) for all \(\gamma \in B\) . We distinguish three cases.

Case 1: \(\gamma \neq \alpha\) and \(\gamma \neq \beta\) . In this case, \(x_{\gamma}\) commutes with \(x_{-\alpha}\) and \(x_{-\beta}\) , and hence with \(y_{\alpha \beta}\) .

Case 2: \(\gamma = \beta\) . In this case, \(x_{\gamma}\) commutes with \(x_{-\alpha}\) , so

\[
\begin{array}{r l} & {[ x _ {\gamma}, y _ {\alpha \beta} ] = (\mathrm{ad} x _ {- \alpha}) ^ {1 - n (\beta , \alpha)} [ x _ {\gamma}, x _ {- \beta} ]} \\ & {\qquad = - (\mathrm{ad} x _ {- \alpha}) ^ {1 - n (\beta , \alpha)} h _ {\beta} = - n (\alpha , \beta) (\mathrm{ad} x _ {- \alpha}) ^ {- n (\beta , \alpha)} x _ {- \alpha}.} \end{array}
\]

If \(n(\beta, \alpha) < 0\) , this expression is zero since (ad \(x_{-\alpha}\) ). \(x_{-\alpha} = 0\) . If \(n(\beta, \alpha) = 0\) , then \(n(\alpha, \beta) = 0\) . In both cases, \([x_{\gamma}, y_{\alpha\beta}] = 0\) .

Case 3: \(\gamma = \alpha\) . In the algebra of endomorphisms of \(\mathfrak{a}\) ,

\[
[ - \operatorname{ad} h _ {\alpha}, \operatorname{ad} x _ {- \alpha} ] = 2 \operatorname{ad} x _ {- \alpha}
\]

and {[}ad \(x_{\alpha}\) , ad \(x_{-\alpha}\) {]} = -ad \(h_{\alpha}\) ; thus, by §1, Lemma 1,

\[
[ \operatorname{ad} x _ {\alpha}, (\operatorname{ad} x _ {- \alpha}) ^ {1 - n (\beta , \alpha)} ] = (1 - n (\beta , \alpha)) (\operatorname{ad} x _ {- \alpha}) ^ {- n (\beta , \alpha)} (- \operatorname{ad} h _ {\alpha} - n (\beta , \alpha)).
\]

Consequently,

\[
\begin{array}{c} [ x _ {\gamma}, y _ {\alpha \beta} ] = [ \text {ad} x _ {\alpha}, (\text {ad} x _ {- \alpha}) ^ {1 - n (\beta , \alpha)} ] x _ {- \beta} + (\text {ad} x _ {- \alpha}) ^ {1 - n (\beta , \alpha)} (\text {ad} x _ {\alpha}) x _ {- \beta} \\ = - (1 - n (\beta , \alpha)) (\text {ad} x _ {- \alpha}) ^ {- n (\beta , \alpha)} (\text {ad} h _ {\alpha} + n (\beta , \alpha)) x _ {- \beta} \\ + (\text {ad} x _ {- \alpha}) ^ {1 - n (\beta , \alpha)} (\text {ad} x _ {\alpha}) x _ {- \beta}. \end{array}
\]

Now \([h_{\alpha},x_{-\beta}] + n(\beta ,\alpha)x_{-\beta} = 0\) and \([x_{\alpha},x_{-\beta}] = 0\) , so \([x_{\gamma},y_{\alpha \beta}] = 0\) .

Lemma 7. The ideal \(\mathfrak{n}\) of \(\mathfrak{a}_{+}\) generated by the \(x_{\alpha \beta} (\alpha, \beta \in \mathrm{B}, \alpha \neq \beta)\) is an ideal of \(\mathfrak{a}\) . The ideal of \(\mathfrak{a}_{-}\) generated by the \(y_{\alpha \beta} (\alpha, \beta \in \mathrm{B}, \alpha \neq \beta)\) is an ideal of \(\mathfrak{a}\) and is equal to \(\theta(\mathfrak{n})\) .

Let \(\mathfrak{n}' = \sum_{\alpha, \beta \in \mathrm{B}, \alpha \neq \beta} kx_{\alpha \beta}\) . Since each \(x_{\alpha \beta}\) is homogeneous in \(\mathfrak{a}\) , \([\mathfrak{a}^0, \mathfrak{n}'] \subset \mathfrak{n}'\) (Lemma 5 and Prop. 2). Let U (resp. V) be the enveloping algebra of \(\mathfrak{a}\) (resp. \(\mathfrak{a}_+\) ), and \(\sigma\) the representation of U on \(\mathfrak{a}\) defined by the adjoint representation of \(\mathfrak{a}\) . The ideal of \(\mathfrak{a}\) generated by \(\mathfrak{n}'\) is \(\sigma(\mathrm{U})\mathfrak{n}'\) . Now \(\mathfrak{a} = \mathfrak{a}_+ + \mathfrak{a}^0 + \mathfrak{a}_-\) (Prop. 2), \(\sigma(\mathfrak{a}_-)\mathfrak{n}' = 0\) (Lemma 6), and \(\sigma(\mathfrak{a}^0)\mathfrak{n}' \subset \mathfrak{n}'\) by the preceding. By the Poincaré-Birkhoff-Witt theorem, \(\sigma(\mathrm{U})\mathfrak{n}' = \sigma(\mathrm{V})\mathfrak{n}'\) , which proves the first assertion of the lemma. It follows that the ideal of \(\theta(\mathfrak{a}_+) = \mathfrak{a}_-\) generated by the \(\theta(x_{\alpha \beta}) = y_{\alpha \beta}\) ( \(\alpha, \beta \in \mathrm{B}, \alpha \neq \beta\) ) is the ideal \(\theta(\mathfrak{n})\) of \(\mathfrak{a}\) . Q.E.D.

The ideal \(\mathfrak{n} + \theta(\mathfrak{n})\) of a is graded since it is generated by homogeneous elements. Consequently, the Lie algebra \(\mathfrak{a}/(\mathfrak{n}+\theta(\mathfrak{n}))\) is a Q-graded Lie algebra; in the remainder of this paragraph, it is denoted by \(g_{B}\) , or simply by g. By Prop. 2, if \(g^{\mu} \neq 0\) then \(\mu \in Q_{+}\) , or \(\mu \in Q_{-}\) , or \(\mu = 0\) . Denote by \(X_{\alpha}\) (resp. \(H_{\alpha}, X_{-\alpha}\) ) the canonical image of \(x_{\alpha}\) (resp. \(h_{\alpha}, x_{-\alpha}\) ) in g. In view of the definition of a, n and \(\theta(\mathfrak{n})\) , it follows that g is the Lie algebra defined by the generating family \(((X_{\alpha}, H_{\alpha}, X_{-\alpha}))_{\alpha \in B}\) and the relations

\[
[ H _ {\alpha}, H _ {\beta} ] = 0\tag{16}
\]

\[
[ H _ {\alpha}, X _ {\beta} ] - n (\beta , \alpha) X _ {\beta} = 0\tag{17}
\]

\[
[ H _ {\alpha}, X _ {- \beta} ] + n (\beta , \alpha) X _ {- \beta} = 0\tag{18}
\]

\[
[ X _ {\alpha}, X _ {- \alpha} ] + H _ {\alpha} = 0\tag{19}
\]

\[
[ X _ {\alpha}, X _ {- \beta} ] = 0 (\alpha \neq \beta)\tag{20}
\]

\[
(\operatorname{ad} X _ {\alpha}) ^ {1 - n (\beta , \alpha)} X _ {\beta} = 0 \quad (\alpha \neq \beta)\tag{21}
\]

\[
(\operatorname{ad} X _ {- \alpha}) ^ {1 - n (\beta , \alpha)} X _ {- \beta} = 0 \quad (\alpha \neq \beta).\tag{22}
\]

Let \(z \in \mathfrak{g}\) and \(\mu \in \mathrm{Q}\) . Then \(z \in \mathfrak{g}^{\mu}\) if and only if \([H_{\alpha}, z] = \langle \mu, \alpha^{\vee} \rangle z\) for all \(\alpha \in \mathrm{B}\) . This follows from Lemma 5.

Since \(\mathfrak{a}^0\cap (\mathfrak{n} + \theta (\mathfrak{n})) = 0\) , the canonical map from \(\mathfrak{a}^0\) to \(\mathfrak{g}^0\) is an isomorphism. Consequently, \((H_{\alpha})_{\alpha \in \mathrm{B}}\) is a basis of the vector space \(\mathfrak{g}^0\) . The commutative subalgebra \(\mathfrak{g}^0\) of \(\mathfrak{g}_{\mathrm{B}}\) will be denoted by \(\mathfrak{h}_{\mathrm{B}}\) or simply by \(\mathfrak{h}\) . There exists a unique isomorphism \(\mu \mapsto \mu_{\mathrm{B}}\) from V to \(\mathfrak{h}^*\) such that \(\langle \mu_{\mathrm{B}},H_{\alpha}\rangle = \langle \mu ,\alpha^{\vee}\rangle\) for all \(\mu \in \mathrm{V}\) and all \(\alpha \in \mathrm{B}\) .

The involutive automorphism \(\theta\) of a defines by passage to the quotient an involutive automorphism of g that will also be denoted by \(\theta\) . We have \(\theta(X_{\alpha}) = X_{-\alpha}\) for \(\alpha \in B \cup (-B)\) , and \(\theta(H_{\alpha}) = -H_{\alpha}\) .

THEOREM 1. Let R be a reduced root system, B a basis of R. Let g be the Lie algebra defined by the generating family \(((X_{\alpha}, H_{\alpha}, X_{-\alpha}))_{\alpha \in \mathrm{B}}\) and the relations (16) to (22). Let \(h = \sum_{\alpha \in B} kH_{\alpha}\) . Then \((\mathfrak{g}, \mathfrak{h})\) is a split semi-simple Lie algebra. The isomorphism \(\mu \mapsto \mu_{B}\) from V to \(h^{*}\) maps R to the root system of \((\mathfrak{g}, \mathfrak{h})\) . For all \(\mu \in R\) , \(g^{\mu}\) is the eigenspace relative to the root \(\mu\) .

The proof follows that of Lemmas 8, 9, 10, 11.

Lemma 8. Let \(\alpha \in \mathrm{B} \cup (-\mathrm{B})\) . Then \(\operatorname{ad} X_{\alpha}\) is locally nilpotent. \footnote{An endomorphism \(u\) of a vector space \(V\) is called locally nilpotent (or almost nilpotent) if, for every \(v \in V\), there exists a positive integer \(n\) such that \(u^n(v) = 0\) (cf. Commutative Algebra, Chap. IV, §1, no. 4, Def. 2). Then \(\exp(u)\), or \(e^u\), is defined by the formula \(e^u(v) = \sum_{n \geq 0} (1/n!)u^n(v)\) for all \(v \in V\).}

Assume that \(\alpha \in \mathrm{B}\) . Let \(\mathfrak{g}'\) be the set of \(z \in \mathfrak{g}\) such that \((\operatorname{ad} X_{\alpha})^p z = 0\) for sufficiently large \(p\) . Since \(\operatorname{ad} X_{\alpha}\) is a derivation of \(\mathfrak{g}\) , \(\mathfrak{g}'\) is a subalgebra of \(\mathfrak{g}\) . By (21), \(X_{\beta} \in \mathfrak{g}'\) for all \(\beta \in \mathrm{B}\) . By (17), (19), (20), \(H_{\beta} \in \mathfrak{g}'\) and \(X_{-\beta} \in \mathfrak{g}'\) for all \(\beta \in \mathrm{B}\) . Hence \(\mathfrak{g}' = \mathfrak{g}\) and \(\operatorname{ad} X_{\alpha}\) is locally nilpotent. Since \(\operatorname{ad} X_{-\alpha} = \theta (\operatorname{ad} X_{\alpha})\theta^{-1}\) , we see that \(\operatorname{ad} X_{-\alpha}\) is locally nilpotent.

We shall see that g is finite dimensional, so that \(ad X_{\alpha}\) is actually nilpotent.

Lemma 9. Let \(\mu, \nu \in \mathbf{Q}\) and \(w \in \mathrm{W}(\mathrm{R})\) be such that \(w\mu = \nu\) . There exists an automorphism of \(\mathfrak{g}\) that takes \(\mathfrak{g}^{\mu}\) to \(\mathfrak{g}^{\nu}\) .

For all \(\alpha \in \mathrm{B}\) , let \(s_{\alpha}\) be the reflection in V defined by \(\alpha\) . Since \(\mathrm{W}(\mathrm{R})\) is generated by the \(s_{\alpha}\) (Chap. VI, §1, no. 5, Remark 1), it suffices to prove the lemma when \(w = s_{\alpha}\) . In view of Lemma 8, we can define

\[
\theta_ {\alpha} = e ^ {\mathrm{ad} X _ {\alpha}} e ^ {\mathrm{ad} X _ {- \alpha}} e ^ {\mathrm{ad} X _ {\alpha}}.
\]

It is verified as in Chap. I, §6, no. 8, that \(\theta_{\alpha}\) is an automorphism of \(\mathfrak{g}\) . We have

\[
\begin{array}{l} \theta_ {\alpha} (H _ {\alpha}) = e ^ {\mathrm{ad} X _ {\alpha}} e ^ {\mathrm{ad} X _ {- \alpha}} (H _ {\beta} - n (\alpha , \beta) X _ {\alpha}) \\ \qquad = e ^ {\mathrm{ad} X _ {\alpha}} \bigg (H _ {\beta} - n (\alpha , \beta) X _ {\alpha} + n (\alpha , \beta) X _ {- \alpha} - n (\alpha , \beta) H _ {\alpha} - \frac {n (\alpha , \beta)}{2} 2 X _ {- \alpha} \bigg) \\ \qquad = e ^ {\mathrm{ad} X _ {\alpha}} \left(H _ {\beta} - n (\alpha , \beta) H _ {\alpha} - n (\alpha , \beta) X _ {\alpha}\right) \\ \qquad = H _ {\beta} - n (\alpha , \beta) H _ {\alpha} - n (\alpha , \beta) X _ {\alpha} - n (\alpha , \beta) X _ {\alpha} - n (\alpha , \beta) (- 2 X _ {\alpha}) \\ \qquad = H _ {\beta} - n (\alpha , \beta) H _ {\alpha}. \end{array}
\]

If \(z\in \mathfrak{g}^{\mu}\)

\[
\begin{array}{r l} & {[ H _ {\beta}, \theta_ {\alpha} ^ {- 1} z ] = \theta_ {\alpha} ^ {- 1} [ H _ {\beta} - n (\alpha , \beta) H _ {\alpha}, z ]} \\ & {\qquad = \theta_ {\alpha} ^ {- 1} (\langle \mu , \beta^ {\vee} \rangle z - n (\alpha , \beta) \langle \mu , \alpha^ {\vee} \rangle z)} \\ & {\qquad = \langle \mu - \langle \alpha^ {\vee}, \mu \rangle \alpha , \beta^ {\vee} \rangle \theta_ {\alpha} ^ {- 1} z = \langle s _ {\alpha} \mu , \beta^ {\vee} \rangle \theta_ {\alpha} ^ {- 1} z,} \end{array}
\]

so \(\theta_{\alpha}^{-1}z\in\mathfrak{g}^{s_{\alpha}\mu}\) . This shows that \(\theta_{\alpha}^{-1}\mathfrak{g}^{\mu}\subset\mathfrak{g}^{s_{\alpha}\mu}\) . Since \(\theta_{\alpha}\) is an automorphism and since this inclusion holds for all \(\mu\in Q\) , we see that \(\theta_{\alpha}^{-1}\mathfrak{g}^{\mu}=\mathfrak{g}^{s_{\alpha}\mu}\) , which proves the lemma.

Lemma 10. Let \(\mu\in Q\) , and assume that \(\mu\) is not a multiple of a root. There exists \(w\in\mathrm{W}(\mathrm{R})\) such that certain of the coordinates of \(w\mu\) with respect to the basis B are \(> 0\) and certain of them are \(< 0\).

Let \(V_{R}\) be the vector space \(Q \otimes_{Z} R\) , in which R is a root system. By the assumption, there exists \(f \in V_{R}^{*}\) such that \(\langle f, \alpha \rangle \neq 0\) for all \(\alpha \in R\) , and \(\langle f, \mu \rangle = 0\) . There exists a chamber C of \(R^{\vee}\) such that \(f \in C\) . By Chap. VI, §1, no. 5, Th. 2 (i), there exists \(w \in W(R)\) such that wf belongs to the chamber associated to B, in other words such that \(\langle wf, \alpha \rangle > 0\) for all \(\alpha \in B\) . Write \(w\mu = \sum_{\alpha \in B} t_{\alpha}\alpha\) . Then

\[
0 = \langle f, \mu \rangle = \langle w f, w \mu \rangle = \sum_ {\alpha \in B} t _ {\alpha} \langle w f, \alpha \rangle ,
\]

which proves that certain \(t_{\alpha}\) are \(> 0\) and others are \(< 0\) .

Lemma 11. Let \(\mu \in \mathrm{Q}\) . If \(\mu \notin \mathrm{R} \cup \{0\}\) , then \(\mathfrak{g}^{\mu} = 0\) . If \(\mu \in \mathrm{R}\) , then \(\dim \mathfrak{g}^{\mu} = 1\) . 1) If \(\mu\) is not a multiple of an element of \(\mathrm{R}\) , there exists \(w \in \mathrm{W}\) such that \(w\mu \notin \mathrm{Q}_{+} \cup \mathrm{Q}_{-}\) (Lemma 10), so \(\mathfrak{a}^{w\mu} = 0\) , \(\mathfrak{g}^{w\mu} = 0\) , and hence \(\mathfrak{g}^{\mu} = 0\) (Lemma 9).

\begin{enumerate}
\def\labelenumi{\arabic{enumi})}
\setcounter{enumi}{1}
\item
  Let \(\alpha \in \mathrm{B}\) and let \(m\) be an integer. Since \(\mathfrak{a}_{+}\) is a free Lie algebra with basic family \((x_{\alpha})_{\alpha \in \mathrm{B}}\) , we have \(\dim \mathfrak{a}^{\alpha} = 1\) and \(\mathfrak{a}^{m\alpha} = 0\) for \(m > 1\) (Chap. II, §2, no. 6, Prop. 4). Hence \(\dim \mathfrak{g}^{\alpha} \leq 1\) and \(\mathfrak{g}^{m\alpha} = 0\) for \(m > 1\) . We cannot have \(\mathfrak{g}^{\alpha} = 0\) , as this would imply that \(x_{\alpha} \in \mathfrak{n} + \theta \mathfrak{n}\) , and hence that \(\mathfrak{n} + \theta \mathfrak{n}\) contains \(h_{\alpha} = -[x_{\alpha}, x_{-\alpha}]\) , whereas \(\mathfrak{a}^0 \cap (\mathfrak{n} + \theta \mathfrak{n}) = 0\) . Consequently, \(\dim \mathfrak{g}^{\alpha} = 1\) .
\item
  If \(\mu \in \mathrm{R}\) , there exists \(w \in \mathrm{W}(\mathrm{R})\) such that \(w(\mu) \in \mathrm{B}\) (Chap. VI, §1, no. 5, Prop. 15), so \(\dim \mathfrak{g}^{\mu} = \dim \mathfrak{g}^{w\mu} = 1\) . Moreover, if \(n\) is an integer \(>1\) then \(\mathfrak{g}^{nw(\mu)} = 0\) and so \(\mathfrak{g}^{n\mu} = 0\) .
\end{enumerate}

Proof of Theorem 1.

Since \(\dim \mathfrak{g}^0 = \operatorname{CardB}\) , it follows from Lemma 11 that \(\mathfrak{g}\) is of finite dimension equal to \(\operatorname{CardB} + \operatorname{CardR}\) . We show that \(\mathfrak{g}\) is semi-simple. Let \(\mathfrak{k}\) be a commutative ideal of \(\mathfrak{g}\) . Since \(\mathfrak{k}\) is stable under \(\operatorname{ad}(\mathfrak{h})\) , \(\mathfrak{k} = (\mathfrak{k} \cap \mathfrak{h}) + \sum_{\mu \in \mathrm{R}} (\mathfrak{k} \cap \mathfrak{g}^\mu)\) .

It is clear that, for all \(\alpha \in \mathrm{B}\) , \(\mathfrak{g}^{\alpha} + \mathfrak{g}^{-\alpha} + kH_{\alpha}\) is isomorphic to \(\mathfrak{sl}(2,k)\) . In view of Lemma 9, for all \(\mu \in \mathrm{R}\) , \(\mathfrak{g}^{\mu}\) is contained in a subalgebra of \(\mathfrak{g}\) isomorphic to \(\mathfrak{sl}(2,k)\) ; consequently, \(\mathfrak{k} \cap \mathfrak{g}^{\mu} = 0\) , so \(\mathfrak{k} \subset \mathfrak{h}\) ; hence

\[
[ \mathfrak {k}, \mathfrak {g} ^ {\mu} ] \subset \mathfrak {k} \cap \mathfrak {g} ^ {\mu} = 0,
\]

so \(\mu_{\mathrm{B}}(\mathfrak{k}) = 0\) for all \(\mu \in \mathrm{R}\) . It follows that \(\mathfrak{k} = 0\) , which proves that \(\mathfrak{g}\) is semi-simple.

Let \(\mu \in \mathrm{R}\) . There exists \(\alpha \in \mathrm{B}\) such that \(\langle \mu, \alpha^{\vee} \rangle \neq 0\) , and \((\operatorname{ad} H_{\alpha})|_{\mathfrak{g}^{\mu}}\) is then a non-zero homothety. Consequently, \(\mathfrak{h}\) is equal to its own normalizer in \(\mathfrak{g}\) , and hence is a Cartan subalgebra of \(\mathfrak{g}\) . For all \(u \in \mathfrak{h}\) , \(\operatorname{ad} u\) is diagonalizable, so \((\mathfrak{g}, \mathfrak{h})\) is a split semi-simple Lie algebra.

For all \(\mu \in \mathrm{R}\) , it is clear that \(\mu_{\mathrm{B}}\) is a root of \((\mathfrak{g},\mathfrak{h})\) and that \(\mathfrak{g}^{\mu}\) is the corresponding eigenspace. The number of roots of \((\mathfrak{g},\mathfrak{h})\) is \(\dim \mathfrak{g} - \dim \mathfrak{h} = \operatorname{Card}\mathrm{R}\) . Hence, the map \(\mu \mapsto \mu_{\mathrm{B}}\) from V to \(\mathfrak{h}^*\) maps R to the root system of \((\mathfrak{g},\mathfrak{h})\) .

